\section{Introduction}\label{chap:introduction}
% Support Vector Machines (SVMs) \cite{cortes1995support} are a cornerstone of convex learning theory.  
% Given labelled data $(x_i,y_i)_{i=1}^M\subset \mathbb{R}^d\times\{-1,+1\}$, an SVM seeks the decision function $F(x)=\operatorname{sgn}(\langle w,x\rangle+b)$ that maximises the geometric margin while tolerating controlled violations via slack variables.  
% Because the resulting optimisation problem is convex, advanced techniques from numerical optimisation—rather than statistical learning alone—determine practical performance.  
% In this project we therefore concentrate on algorithmic aspects: deriving the dual quadratic programme, analysing its properties, and designing scalable solvers.

Support Vector Machines (SVMs) are a theoretically grounded method 
for binary classification, particularly effective in high-dimensional 
settings. Given a set of labelled training data, the aim is to 
construct a decision function that separates the two classes with 
maximal margin while tolerating potential misclassifications. This 
leads to the formulation of the soft-margin SVM, a convex 
optimization problem balancing model complexity and classification 
error via a regularization parameter.

This report presents a theoretical and numerical investigation of 
the soft-margin SVM formulation. The existence and uniqueness of 
solutions is discussed through exploration of the structure and 
behavior of the optimization landscape using convex analysis.
A numerical method is implemented for solving soft-margin SVMs,
and tested on numerical experiments demonstrating the performance of 
the algorithm.

